\clearpage
\scanpage{80}
% Short oblique impression at the upper left of the sheet.
\reversemarginpar
\margin{\figure{\tikz[x=1em,y=1em]{\draw[line width=.23pt] (0,0) -- (1.5,1.25); \fill (0,0) -- (.15,.02) -- (.25,.19) -- (.14,.19) -- cycle; \fill (1.37,1.11) -- (1.47,1.27) -- (1.53,1.22) -- (1.43,1.13) -- cycle;}}}
\normalmarginpar
\mspage{X-1}

\begin{center}
List of elliptic curves with complex multiplication\\
whose modular invariant $\mathrm{j}$ is rational

\medskip
-{}-{}-
\end{center}

These curves are obtained from the orders of imaginary quadratic fields of class number $1$. One must therefore :

a) make the list of imaginary quadratic fields $\underline{\underline{\mathrm{Q}}}(\handwritten{\sqrt}{-\mathrm{d}})$ whose maximal order (i.e. the ring of integers) is principal. The result is :
\reversemarginpar
\margin{\handwritten{\uline{NB}\quad the $\mathrm{d}$\\are prime !}}
\normalmarginpar
\[
\begin{aligned}
&\mathrm{d} = \handwritten{\overbrace}{1,\,2,\,3,\,7,\,11}^{\inserted{\handwritten{\text{``\unsure{Euclidean|Euclidean}''}}}},\,19,\,43,\,67,\,163 \quad \text{plus perhaps}\\
&\text{one other value}\quad \mathrm{d}_{?}\quad .
\end{aligned}
\]

b) For each of the values of $\mathrm{d}$ above, determine which conductors $\underline{\mathrm{f}}$ lead to an order having class number equal to $1$.

One finds $\underline{\mathrm{f}}=1,2$ for $\mathrm{d}=1$ , $\underline{\mathrm{f}}=1,2,3$ for $\mathrm{d}=3$, $\underline{\mathrm{f}}=1,2$ for $\mathrm{d}=7$, and $\underline{\mathrm{f}}=1$ for $\mathrm{d}=2,\,11,\,19,\,43,\,67,\,163$ and $\mathrm{d}_{?}$ (if it exists).

c) For each pair $(\mathrm{d},\underline{\mathrm{f}})$, take a basis $(1,\mathrm{t})$ of the corresponding order, and calculate $\mathrm{j}(\mathrm{t})$; this is the modular invariant sought; it is rational.

The values of $\mathrm{j}(\mathrm{t})$ are all found in H.WEBER, \uline{Lehrbuch der Algebra}, volume III; more precisely, Weber gives the values of two functions $\mathrm{f}(\mathrm{t})$ and $\mathrm{f}_{1}(\mathrm{t})$ related to $\mathrm{j}$ by the relations
\[
\mathrm{j}=(\mathrm{f}^{24}-16)^{3}/\mathrm{f}^{24}
       =(\mathrm{f}_{1}^{24}+16)^{3}/\mathrm{f}_{1}^{24}\quad .
\]

One thus obtains the following results :

\begingroup
\makeatletter\tagsleft@true\makeatother
\[
(\mathrm{d}=1,\ \underline{\mathrm{f}}=1)\quad \rightarrow\quad
\mathrm{j}=\mathrm{j}(\mathrm{i})=2^{6}3^{3}
\tag*{1.1.}
\]
(elliptic curve having an automorphism group of order $4$)

\[
(\mathrm{d}=1,\ \underline{\mathrm{f}}=2)\quad \rightarrow\quad
\mathrm{j}=\mathrm{j}(2\mathrm{i})=(2.3.11)^{3}
\tag*{1.2}
\]
\endgroup
% archive sheet number: 79
